% ID: SEL-D0022-OA-E04
% ESTADO: validada. S023, revisada con los criterios oficiales.
\textbf{a)} El centro del intervalo es la media muestral:
$\bar x=(164.86+171.14)/2=\boxed{168\text{ cm}}$.
La semiamplitud impresa es $3.14$ cm y el nuevo margen debe ser
como máximo $1.57$ cm. Manteniendo la confianza al $95\%$,
\[
E=\frac{z_{0.975}\,8}{\sqrt n}\le1.57
\quad\Longleftrightarrow\quad
n\ge\left(\frac{1.959964\cdot8}{1.57}\right)^2\approx99.74.
\]
El mínimo es $\boxed{100}$ individuos. Como el margen es inversamente
proporcional a $\sqrt n$, reducirlo a la mitad equivale a cuadruplicar
el tamaño muestral si se toma como referencia el margen sin redondear.

\textbf{b)} Suponemos \emph{afijación proporcional}: la fracción
de muestreo es igual en cada estrato. Como $5+7+15=27$ y
$243/27=9$, se multiplica cada recuento por 9:
45 en Yoga, 63 en Pilates y 135 en Mantenimiento.
La suma es 243. Con otra afijación, el muestreo estratificado por sí
solo no determina los tamaños de los estratos.
